\section{積分区間の変換}\label{節:置換と積分区間}

不定積分では最後に $x$ へ戻したが,定積分ならば新しい変数のまま計算を終えられる.根拠は\sectionref{節:合成形の置換}の定理の後半
\[
 \int_a^b f(g(x))g'(x)\dd x
 =\int_{g(a)}^{g(b)}f(t)\dd t
\]
である.\,$t=g(x)$ と置いたら,\,$x=a$ に対応する $t=g(a)$ を下端に,\,$x=b$ に対応する $t=g(b)$ を上端にする.大きい数を自動的に上端にするのではなく,元の端点との対応を保つ.そのため,計算の初めに $x$ と $t$ の端点の対応を
\[
 \begin{array}{c|ccc}
 x&a&\to&b\\ \hline
 t&g(a)&\to&g(b)
 \end{array}
\]
のような表にして書いておくとよい.

\begin{bookpoint}
\textbf{POINT}\par
$t=g(x)$ と置いたら端点の対応表を書き,\,$g(a)$ を下端,\,$g(b)$ を上端にして,\,$x$ に戻さずに $t$ のまま計算する.
\end{bookpoint}

\begin{example}{途中で折り返す置換}
$\int_{-1}^{2}2x\cos(x^2)\dd x$ を求めよ.
\end{example}

\paragraphbf{方針}
$t=x^2$ と置くと $\dd t=2x\dd x$ である.端点 $x=-1,2$ を $t=1,4$ に移して,\,$t$ のまま計算する.

\paragraphbf{解答}
$\sin4-\sin1$

\paragraphbf{解法}
$t=x^2$ と置けば $\dd t=2x\dd x$ であり,端点の対応は
\[
 \begin{array}{c|ccc}
 x&-1&\to&2\\ \hline
 t&1&\to&4
 \end{array}
\]
である.したがって合成形の置換積分から
\begin{align*}
 \int_{-1}^{2}2x\cos(x^2)\dd x
 &=\int_1^4\cos t\dd t\\
 &=\sin4-\sin1
\end{align*}
を得る.角 $1,4$ はラジアンである.\,$-1\le x\le2$ で $t=x^2$ は単調ではないが,定理に単調性の条件はないので,端点だけを見たこの計算で正しい.

\paragraphbf{折り返しの確認}
端点だけを見た式が正しいのは,途中の折り返しが定積分の向きに反映されているためである.図\ref{図:置換の折返し}に対応を示す.\,$x$ が $-1$ から $0$ まで増える間に $t$ は $1$ から $0$ へ減り,その後 $4$ へ増える.\,$x=0$ で積分を分けて確かめると,
\begin{align*}
 \int_{-1}^{2}2x\cos(x^2)\dd x
 &=\int_1^0\cos t\dd t+\int_0^4\cos t\dd t\\
 &=-\int_0^1\cos t\dd t+\int_0^4\cos t\dd t\\
 &=\int_1^4\cos t\dd t
\end{align*}
となる.往復した範囲 $0\le t\le1$ の寄与が符号付きで打ち消し合っているのだ.この計算は,曲線の下の面積を重複も含めて足す操作とは違う.

\begin{figure}[htbp]
 \centering
 \begin{tikzpicture}[x=1.6cm,y=0.8cm]
  \draw[-{Stealth}] (-1.35,0)--(2.35,0) node[right] {$x$};
  \draw[-{Stealth}] (0,-0.35)--(0,4.65) node[above] {$t$};
  \draw[thick,blue,domain=-1:2,samples=65] plot (\x,{\x*\x});
  \draw[dashed] (-1,0)--(-1,1)--(0,1);
  \draw[dashed] (2,0)--(2,4)--(0,4);
  \fill (-1,1) circle (1.5pt);
  \fill (2,4) circle (1.5pt);
  \node[below] at (-1,0) {$-1$};
  \node[below right] at (0,0) {$0$};
  \node[below] at (2,0) {$2$};
  \node[left] at (0,1) {$1$};
  \node[left] at (0,4) {$4$};
  \node[right] at (1.6,1.9) {$t=x^2$};
  \draw[-{Stealth},thick] (-0.88,0.85)--(-0.55,0.41);
  \draw[-{Stealth},thick] (1.12,1.38)--(1.43,2.17);
 \end{tikzpicture}
 \caption{$x$ は増え続けるが,\,$t=x^2$ は減少してから増加する.}
 \label{図:置換の折返し}
\end{figure}

\paragraphbf{減少する置換と2つの負号}
$g$ が減少する置換では,負号が2か所に現れる.たとえば $t=1-x$ と置くと,\,$x$ が $0$ から $1$ へ動くとき,\,$t$ は $1$ から $0$ へ動く.さらに $\dd t=-\dd x$ だから,
\begin{align*}
 \int_0^1(1-x)^2\dd x
 &=-\int_1^0 t^2\dd t\\
 &=\int_0^1t^2\dd t\\
 &=\frac13
\end{align*}
となる.1行目の負号は微分 $\dd t=-\dd x$ から,2行目で負号が消えたのは上下端を入れ換えたからであり,それぞれ別の理由で現れる.途中で片方だけを落とさないよう,対応表と $\dd t$ を先に書いておくとよい.
